% Generated draft: requires mathematical and source review.
\documentclass[11pt]{article}
\usepackage[margin=25mm]{geometry}
\usepackage{fontspec}
\setmainfont{Latin Modern Roman}
\usepackage{amsmath,amssymb,mathtools,unicode-math}
\setmathfont{Latin Modern Math}
\usepackage{enumitem,needspace,xurl,hyperref}
\setlength{\parindent}{0pt}
\setlength{\parskip}{5pt}
\setlength{\emergencystretch}{3em}
\newcommand{\N}{\mathbb{N}}
\newcommand{\Z}{\mathbb{Z}}
\newcommand{\Q}{\mathbb{Q}}
\newcommand{\R}{\mathbb{R}}
\newcommand{\C}{\mathbb{C}}
\newcommand{\F}{\mathbb{F}}
\newcommand{\A}{\mathbb{A}}
\newcommand{\PP}{\mathbb{P}}
\newcommand{\E}{\mathbb{E}}
\newcommand{\eps}{\varepsilon}
\newcommand{\dd}{\,\mathrm{d}}
\begin{document}
\section*{UnsolvedMath ID: 701110}
\textbf{Source problem code:} catalogue entry
\subsection{Definitions and mathematical statement}
For a set \(\Omega\), let \(\operatorname{Sym}(\Omega)\) denote its group of permutations. For \(g\in\operatorname{Sym}(\Omega)\), define its support by
\[
\operatorname{supp}_{\Omega}(g)=\{\omega\in\Omega:g(\omega)\ne\omega\}.
\]
A permutation group \(G\leq\operatorname{Sym}(\Omega)\) is \emph{finitary} if \(\operatorname{supp}_{\Omega}(g)\) is finite for every \(g\in G\). Its action is transitive if, for every \(\alpha,\beta\in\Omega\), some \(g\in G\) satisfies \(g(\alpha)=\beta\).

The supplied question asks whether there exist finitary barely transitive permutation groups. The record does not define \emph{barely transitive} or explicitly specify the cardinality of the permutation domain, so a fully determined formulation requires clarification of that term.

For precision, the following is a conditional formulation only: if \emph{barely transitive} means a transitive action on an infinite set in which every orbit of every proper subgroup is finite, the question becomes whether there exist an infinite set \(\Omega\) and \(G\leq\operatorname{Sym}(\Omega)\) such that
\[
\begin{gathered}
\forall\alpha,\beta\in\Omega\;\exists g\in G:\quad g(\alpha)=\beta,\\
\forall g\in G:\quad |\operatorname{supp}_{\Omega}(g)|<\infty,\\
\forall H<G\;\forall\omega\in\Omega:\quad |H\omega|<\infty,
\end{gathered}
\]
where \(H<G\) means that \(H\) is a proper subgroup and \(H\omega=\{h(\omega):h\in H\}\). The supplied record does not establish this interpretation of barely transitive; the infinite-domain requirement and the proper-subgroup orbit condition require confirmation from the cited source.
\subsection{Short English statement}
V. V. Belyaev asks whether a barely transitive permutation group can have every element move only finitely many points. The source record leaves barely transitive undefined. The displayed formulation is conditional on interpreting it as transitivity on an infinite set with finite orbits for every proper subgroup.
\subsection{Sources}
\begin{itemize}
\item Local UnsolvedMath registry, permanent ID \(701110\), exact source problem code \texttt{catalogue entry}, title \emph{Kourovka Notebook Problem 15.22}; supplied upstream revision \texttt{d371ee917ac9647fbb06cd6fc529deec6f3dda82}.
\item Supplied source attribution: \emph{The Kourovka Notebook (Unsolved Problems in Group Theory)}, original Problem 15.22, 15th issue (2002), proposed by V. V. Belyaev. The background also identifies the 21st issue edition (2026), \texttt{arXiv:1401.0300v48}.
\item The supplied collection description names E. I. Khukhro and V. D. Mazurov as editors. The registry labels this problem \emph{open}; this is supplied metadata, not an independently verified mathematical or literature assessment. The cited source has not been independently consulted for this draft.
\end{itemize}
\end{document}
